\cotefile{11}
\scanpage{16}
\mspage{-2.29-}
% archive sheet number: 15

\uline{Proof}: Immediate, since $\overline{\lambda_1\unsure{\circ|\cdot}\lambda_2}=\bar{\lambda}_1+\bar{\lambda}_2$.

\uline{Corollary} \uline{2.16}: Let a reductive $G$ act on $X$, proper over $k$. Let $L\in\operatorname{Pic}^{G}(X)$ be ample, and let $x\in X_k$. Then $\{\delta\mid \unsure{\nu|v}^{L}(x;\delta)\leq 0\}$ is semi-convex, and $\{\delta\mid \nu^{L}(x,\delta)<0\}$ is convex.

\uline{Proof}: Immediate.

It follows that from this Corollary and from Theorem 2.1, that if \deleted{\unsure{$($}}$x\in X_k$ is not \uline{semi-stable}, then we can associate to $x$ the convex set $C$ of $\delta\in\Delta(G)$ for which $\nu^{L}(x,\delta)<0$. Roughly speaking, we can say that a parabolic subgroup $P$ is ``responsible'' for this breakdown of semi-stability if $C\cap\Delta_P(G)\neq\unsure{\phi|\emptyset}$. Since $C$ is convex, not too many parabolic subgroups $P$ are responsible; according to a conjecture of Tits, there is \inserted{\handwritten{even}} a natural way to find one $P$ which is \uline{most} responsible, i.e. $P(\delta)$ for $\delta=$ the ``center'' of $C$ (cf. Preface). Another consequence of Proposition 2.14 is:

\uline{Proposition} \uline{2.17}: Let a reductive $G$ act via $\sigma$ on $X$, proper over $k$. Let $L\in\operatorname{Pic}^{G}(x)$ be ample. Then there is a constant $K$ such that
\[
 \lvert\nu^{L}(x,\delta)\rvert\leq K
\]
for all $x\in X_k$ and $\delta\in\Delta(G)$. Moreover, for every subset $\handwritten{S\deleted{\unsure{\in}}\subset X}_k$, there is a $\delta_{\unsure{0|o}}\in\Delta(G)$ and $x_0\in S$ such that
\[
 \nu^{L}(x_0,\delta_0)\leq\nu^{L}(x,\delta)^{\unsure{\circ|.}}
\]
for every other $\delta\in\Delta(G)$, $x\in S$.
